The diagrams in this section vary one input at a time for Au, Ni, Cu, Pd, Pt, and Ti, at 298.15~K. The reference conditions are a nominal aqueous metal-ion concentration of $10^{-4}\,\mathrm{M}$, $[\ce{NH3}]_{\mathrm{tot}}=0.02\,\mathrm{M}$, $[\mathrm{Gly}]_{\mathrm{tot}}=0.1\,\mathrm{M}$, and no added \ce{CN-} or \ce{NO2-}. The cyanide sweep additionally examines $10^{-4}$ and $10^{-3}\,\mathrm{M}$ \ce{CN-}. Each black outline denotes pH 10.5--13.5 and $-2.0$ to $2.3$~V vs RHE; the plotted potential axis is vs SHE.

% Au sensitivity panels
\begin{figure}[!htbp]
    \centering
    \includegraphics[width=\textwidth]{Au_activity_three_levels_1x3_side_colors.pdf}
    \caption{Au Pourbaix sensitivity to nominal aqueous metal-ion concentration. Panels (a)--(c) use $10^{-6}$, $10^{-4}$, and $10^{-2}\,\mathrm{M}$, respectively; fixed conditions are $[\ce{NH3}]_{\mathrm{tot}}=0.02\,\mathrm{M}$, $[\mathrm{Gly}]_{\mathrm{tot}}=0.1\,\mathrm{M}$, and no added cyanide. The black outline denotes pH 10.5--13.5 and $-2.0$ to $2.3$~V vs RHE (shown on the SHE potential axis).}
    \label{fig:Au_activity_sensitivity_SI}
\end{figure}

\begin{figure}[!htbp]
    \centering
    \includegraphics[width=\textwidth]{Au_NH3_three_levels_1x3_side_colors.pdf}
    \caption{Au Pourbaix sensitivity to total ammonia concentration. Panels (a)--(c) use $0.002$, $0.02$, and $0.1\,\mathrm{M}$, respectively; fixed conditions are nominal metal-ion concentration $10^{-4}\,\mathrm{M}$, $[\mathrm{Gly}]_{\mathrm{tot}}=0.1\,\mathrm{M}$, and no added cyanide. The black outline denotes pH 10.5--13.5 and $-2.0$ to $2.3$~V vs RHE (shown on the SHE potential axis).}
    \label{fig:Au_NH3_sensitivity_SI}
\end{figure}

\begin{figure}[!htbp]
    \centering
    \includegraphics[width=\textwidth]{Au_Gly_three_levels_1x3_side_colors.pdf}
    \caption{Au Pourbaix sensitivity to total glycine concentration. Panels (a)--(c) use $0.01$, $0.1$, and $1.0\,\mathrm{M}$, respectively; fixed conditions are nominal metal-ion concentration $10^{-4}\,\mathrm{M}$, $[\ce{NH3}]_{\mathrm{tot}}=0.02\,\mathrm{M}$, and no added cyanide. The black outline denotes pH 10.5--13.5 and $-2.0$ to $2.3$~V vs RHE (shown on the SHE potential axis).}
    \label{fig:Au_Gly_sensitivity_SI}
\end{figure}

\begin{figure}[!htbp]
    \centering
    \includegraphics[width=\textwidth]{Au_CN_three_levels_1x3_side_colors.pdf}
    \caption{Au Pourbaix sensitivity to total cyanide concentration. Panels (a)--(c) use $0$, $10^{-4}$, and $10^{-3}\,\mathrm{M}$, respectively; fixed conditions are nominal metal-ion concentration $10^{-4}\,\mathrm{M}$, $[\ce{NH3}]_{\mathrm{tot}}=0.02\,\mathrm{M}$, and $[\mathrm{Gly}]_{\mathrm{tot}}=0.1\,\mathrm{M}$. The black outline denotes pH 10.5--13.5 and $-2.0$ to $2.3$~V vs RHE (shown on the SHE potential axis).}
    \label{fig:Au_CN_sensitivity_SI}
\end{figure}

% Ni sensitivity panels
\begin{figure}[!htbp]
    \centering
    \includegraphics[width=\textwidth]{Ni_activity_three_levels_1x3_side_colors.pdf}
    \caption{Ni Pourbaix sensitivity to nominal aqueous metal-ion concentration. Panels (a)--(c) use $10^{-6}$, $10^{-4}$, and $10^{-2}\,\mathrm{M}$, respectively; fixed conditions are $[\ce{NH3}]_{\mathrm{tot}}=0.02\,\mathrm{M}$, $[\mathrm{Gly}]_{\mathrm{tot}}=0.1\,\mathrm{M}$, and no added cyanide. The black outline denotes pH 10.5--13.5 and $-2.0$ to $2.3$~V vs RHE (shown on the SHE potential axis).}
    \label{fig:Ni_activity_sensitivity_SI}
\end{figure}

\begin{figure}[!htbp]
    \centering
    \includegraphics[width=\textwidth]{Ni_NH3_three_levels_1x3_side_colors.pdf}
    \caption{Ni Pourbaix sensitivity to total ammonia concentration. Panels (a)--(c) use $0.002$, $0.02$, and $0.1\,\mathrm{M}$, respectively; fixed conditions are nominal metal-ion concentration $10^{-4}\,\mathrm{M}$, $[\mathrm{Gly}]_{\mathrm{tot}}=0.1\,\mathrm{M}$, and no added cyanide. The black outline denotes pH 10.5--13.5 and $-2.0$ to $2.3$~V vs RHE (shown on the SHE potential axis).}
    \label{fig:Ni_NH3_sensitivity_SI}
\end{figure}

\begin{figure}[!htbp]
    \centering
    \includegraphics[width=\textwidth]{Ni_Gly_three_levels_1x3_side_colors.pdf}
    \caption{Ni Pourbaix sensitivity to total glycine concentration. Panels (a)--(c) use $0.01$, $0.1$, and $1.0\,\mathrm{M}$, respectively; fixed conditions are nominal metal-ion concentration $10^{-4}\,\mathrm{M}$, $[\ce{NH3}]_{\mathrm{tot}}=0.02\,\mathrm{M}$, and no added cyanide. The black outline denotes pH 10.5--13.5 and $-2.0$ to $2.3$~V vs RHE (shown on the SHE potential axis).}
    \label{fig:Ni_Gly_sensitivity_SI}
\end{figure}

\begin{figure}[!htbp]
    \centering
    \includegraphics[width=\textwidth]{Ni_CN_three_levels_1x3_side_colors.pdf}
    \caption{Ni Pourbaix sensitivity to total cyanide concentration. Panels (a)--(c) use $0$, $10^{-4}$, and $10^{-3}\,\mathrm{M}$, respectively; fixed conditions are nominal metal-ion concentration $10^{-4}\,\mathrm{M}$, $[\ce{NH3}]_{\mathrm{tot}}=0.02\,\mathrm{M}$, and $[\mathrm{Gly}]_{\mathrm{tot}}=0.1\,\mathrm{M}$. The black outline denotes pH 10.5--13.5 and $-2.0$ to $2.3$~V vs RHE (shown on the SHE potential axis).}
    \label{fig:Ni_CN_sensitivity_SI}
\end{figure}

% Cu sensitivity panels
\begin{figure}[!htbp]
    \centering
    \includegraphics[width=\textwidth]{Cu_activity_three_levels_1x3_side_colors.pdf}
    \caption{Cu Pourbaix sensitivity to nominal aqueous metal-ion concentration. Panels (a)--(c) use $10^{-6}$, $10^{-4}$, and $10^{-2}\,\mathrm{M}$, respectively; fixed conditions are $[\ce{NH3}]_{\mathrm{tot}}=0.02\,\mathrm{M}$, $[\mathrm{Gly}]_{\mathrm{tot}}=0.1\,\mathrm{M}$, and no added cyanide. The black outline denotes pH 10.5--13.5 and $-2.0$ to $2.3$~V vs RHE (shown on the SHE potential axis).}
    \label{fig:Cu_activity_sensitivity_SI}
\end{figure}

\begin{figure}[!htbp]
    \centering
    \includegraphics[width=\textwidth]{Cu_NH3_three_levels_1x3_side_colors.pdf}
    \caption{Cu Pourbaix sensitivity to total ammonia concentration. Panels (a)--(c) use $0.002$, $0.02$, and $0.1\,\mathrm{M}$, respectively; fixed conditions are nominal metal-ion concentration $10^{-4}\,\mathrm{M}$, $[\mathrm{Gly}]_{\mathrm{tot}}=0.1\,\mathrm{M}$, and no added cyanide. The black outline denotes pH 10.5--13.5 and $-2.0$ to $2.3$~V vs RHE (shown on the SHE potential axis).}
    \label{fig:Cu_NH3_sensitivity_SI}
\end{figure}

\begin{figure}[!htbp]
    \centering
    \includegraphics[width=\textwidth]{Cu_Gly_three_levels_1x3_side_colors.pdf}
    \caption{Cu Pourbaix sensitivity to total glycine concentration. Panels (a)--(c) use $0.01$, $0.1$, and $1.0\,\mathrm{M}$, respectively; fixed conditions are nominal metal-ion concentration $10^{-4}\,\mathrm{M}$, $[\ce{NH3}]_{\mathrm{tot}}=0.02\,\mathrm{M}$, and no added cyanide. The black outline denotes pH 10.5--13.5 and $-2.0$ to $2.3$~V vs RHE (shown on the SHE potential axis).}
    \label{fig:Cu_Gly_sensitivity_SI}
\end{figure}

\begin{figure}[!htbp]
    \centering
    \includegraphics[width=\textwidth]{Cu_CN_three_levels_1x3_side_colors.pdf}
    \caption{Cu Pourbaix sensitivity to total cyanide concentration. Panels (a)--(c) use $0$, $10^{-4}$, and $10^{-3}\,\mathrm{M}$, respectively; fixed conditions are nominal metal-ion concentration $10^{-4}\,\mathrm{M}$, $[\ce{NH3}]_{\mathrm{tot}}=0.02\,\mathrm{M}$, and $[\mathrm{Gly}]_{\mathrm{tot}}=0.1\,\mathrm{M}$. The black outline denotes pH 10.5--13.5 and $-2.0$ to $2.3$~V vs RHE (shown on the SHE potential axis).}
    \label{fig:Cu_CN_sensitivity_SI}
\end{figure}

% Pd sensitivity panels
\begin{figure}[!htbp]
    \centering
    \includegraphics[width=\textwidth]{Pd_activity_three_levels_1x3_side_colors.pdf}
    \caption{Pd Pourbaix sensitivity to nominal aqueous metal-ion concentration. Panels (a)--(c) use $10^{-6}$, $10^{-4}$, and $10^{-2}\,\mathrm{M}$, respectively; fixed conditions are $[\ce{NH3}]_{\mathrm{tot}}=0.02\,\mathrm{M}$, $[\mathrm{Gly}]_{\mathrm{tot}}=0.1\,\mathrm{M}$, and no added cyanide. The black outline denotes pH 10.5--13.5 and $-2.0$ to $2.3$~V vs RHE (shown on the SHE potential axis).}
    \label{fig:Pd_activity_sensitivity_SI}
\end{figure}

\begin{figure}[!htbp]
    \centering
    \includegraphics[width=\textwidth]{Pd_NH3_three_levels_1x3_side_colors.pdf}
    \caption{Pd Pourbaix sensitivity to total ammonia concentration. Panels (a)--(c) use $0.002$, $0.02$, and $0.1\,\mathrm{M}$, respectively; fixed conditions are nominal metal-ion concentration $10^{-4}\,\mathrm{M}$, $[\mathrm{Gly}]_{\mathrm{tot}}=0.1\,\mathrm{M}$, and no added cyanide. The black outline denotes pH 10.5--13.5 and $-2.0$ to $2.3$~V vs RHE (shown on the SHE potential axis).}
    \label{fig:Pd_NH3_sensitivity_SI}
\end{figure}

\begin{figure}[!htbp]
    \centering
    \includegraphics[width=\textwidth]{Pd_Gly_three_levels_1x3_side_colors.pdf}
    \caption{Pd Pourbaix sensitivity to total glycine concentration. Panels (a)--(c) use $0.01$, $0.1$, and $1.0\,\mathrm{M}$, respectively; fixed conditions are nominal metal-ion concentration $10^{-4}\,\mathrm{M}$, $[\ce{NH3}]_{\mathrm{tot}}=0.02\,\mathrm{M}$, and no added cyanide. The black outline denotes pH 10.5--13.5 and $-2.0$ to $2.3$~V vs RHE (shown on the SHE potential axis).}
    \label{fig:Pd_Gly_sensitivity_SI}
\end{figure}

\begin{figure}[!htbp]
    \centering
    \includegraphics[width=\textwidth]{Pd_CN_three_levels_1x3_side_colors.pdf}
    \caption{Pd Pourbaix sensitivity to total cyanide concentration. Panels (a)--(c) use $0$, $10^{-4}$, and $10^{-3}\,\mathrm{M}$, respectively; fixed conditions are nominal metal-ion concentration $10^{-4}\,\mathrm{M}$, $[\ce{NH3}]_{\mathrm{tot}}=0.02\,\mathrm{M}$, and $[\mathrm{Gly}]_{\mathrm{tot}}=0.1\,\mathrm{M}$. The black outline denotes pH 10.5--13.5 and $-2.0$ to $2.3$~V vs RHE (shown on the SHE potential axis).}
    \label{fig:Pd_CN_sensitivity_SI}
\end{figure}

% Pt sensitivity panels
\begin{figure}[!htbp]
    \centering
    \includegraphics[width=\textwidth]{Pt_activity_three_levels_1x3_side_colors.pdf}
    \caption{Pt Pourbaix sensitivity to nominal aqueous metal-ion concentration. Panels (a)--(c) use $10^{-6}$, $10^{-4}$, and $10^{-2}\,\mathrm{M}$, respectively; fixed conditions are $[\ce{NH3}]_{\mathrm{tot}}=0.02\,\mathrm{M}$, $[\mathrm{Gly}]_{\mathrm{tot}}=0.1\,\mathrm{M}$, and no added cyanide. The black outline denotes pH 10.5--13.5 and $-2.0$ to $2.3$~V vs RHE (shown on the SHE potential axis).}
    \label{fig:Pt_activity_sensitivity_SI}
\end{figure}

\begin{figure}[!htbp]
    \centering
    \includegraphics[width=\textwidth]{Pt_NH3_three_levels_1x3_side_colors.pdf}
    \caption{Pt Pourbaix sensitivity to total ammonia concentration. Panels (a)--(c) use $0.002$, $0.02$, and $0.1\,\mathrm{M}$, respectively; fixed conditions are nominal metal-ion concentration $10^{-4}\,\mathrm{M}$, $[\mathrm{Gly}]_{\mathrm{tot}}=0.1\,\mathrm{M}$, and no added cyanide. The black outline denotes pH 10.5--13.5 and $-2.0$ to $2.3$~V vs RHE (shown on the SHE potential axis).}
    \label{fig:Pt_NH3_sensitivity_SI}
\end{figure}

\begin{figure}[!htbp]
    \centering
    \includegraphics[width=\textwidth]{Pt_Gly_three_levels_1x3_side_colors.pdf}
    \caption{Pt Pourbaix sensitivity to total glycine concentration. Panels (a)--(c) use $0.01$, $0.1$, and $1.0\,\mathrm{M}$, respectively; fixed conditions are nominal metal-ion concentration $10^{-4}\,\mathrm{M}$, $[\ce{NH3}]_{\mathrm{tot}}=0.02\,\mathrm{M}$, and no added cyanide. The black outline denotes pH 10.5--13.5 and $-2.0$ to $2.3$~V vs RHE (shown on the SHE potential axis).}
    \label{fig:Pt_Gly_sensitivity_SI}
\end{figure}

\begin{figure}[!htbp]
    \centering
    \includegraphics[width=\textwidth]{Pt_CN_three_levels_1x3_side_colors.pdf}
    \caption{Pt Pourbaix sensitivity to total cyanide concentration. Panels (a)--(c) use $0$, $10^{-4}$, and $10^{-3}\,\mathrm{M}$, respectively; fixed conditions are nominal metal-ion concentration $10^{-4}\,\mathrm{M}$, $[\ce{NH3}]_{\mathrm{tot}}=0.02\,\mathrm{M}$, and $[\mathrm{Gly}]_{\mathrm{tot}}=0.1\,\mathrm{M}$. The black outline denotes pH 10.5--13.5 and $-2.0$ to $2.3$~V vs RHE (shown on the SHE potential axis).}
    \label{fig:Pt_CN_sensitivity_SI}
\end{figure}

% Ti sensitivity panels
\begin{figure}[!htbp]
    \centering
    \includegraphics[width=\textwidth]{Ti_activity_three_levels_1x3_side_colors.pdf}
    \caption{Ti Pourbaix sensitivity to nominal aqueous metal-ion concentration. Panels (a)--(c) use $10^{-6}$, $10^{-4}$, and $10^{-2}\,\mathrm{M}$, respectively; fixed conditions are $[\ce{NH3}]_{\mathrm{tot}}=0.02\,\mathrm{M}$, $[\mathrm{Gly}]_{\mathrm{tot}}=0.1\,\mathrm{M}$, and no added cyanide. The black outline denotes pH 10.5--13.5 and $-2.0$ to $2.3$~V vs RHE (shown on the SHE potential axis).}
    \label{fig:Ti_activity_sensitivity_SI}
\end{figure}

\begin{figure}[!htbp]
    \centering
    \includegraphics[width=\textwidth]{Ti_NH3_three_levels_1x3_side_colors.pdf}
    \caption{Ti Pourbaix sensitivity to total ammonia concentration. Panels (a)--(c) use $0.002$, $0.02$, and $0.1\,\mathrm{M}$, respectively; fixed conditions are nominal metal-ion concentration $10^{-4}\,\mathrm{M}$, $[\mathrm{Gly}]_{\mathrm{tot}}=0.1\,\mathrm{M}$, and no added cyanide. The black outline denotes pH 10.5--13.5 and $-2.0$ to $2.3$~V vs RHE (shown on the SHE potential axis).}
    \label{fig:Ti_NH3_sensitivity_SI}
\end{figure}

\begin{figure}[!htbp]
    \centering
    \includegraphics[width=\textwidth]{Ti_Gly_three_levels_1x3_side_colors.pdf}
    \caption{Ti Pourbaix sensitivity to total glycine concentration. Panels (a)--(c) use $0.01$, $0.1$, and $1.0\,\mathrm{M}$, respectively; fixed conditions are nominal metal-ion concentration $10^{-4}\,\mathrm{M}$, $[\ce{NH3}]_{\mathrm{tot}}=0.02\,\mathrm{M}$, and no added cyanide. The black outline denotes pH 10.5--13.5 and $-2.0$ to $2.3$~V vs RHE (shown on the SHE potential axis).}
    \label{fig:Ti_Gly_sensitivity_SI}
\end{figure}

\begin{figure}[!htbp]
    \centering
    \includegraphics[width=\textwidth]{Ti_CN_three_levels_1x3_side_colors.pdf}
    \caption{Ti Pourbaix sensitivity to total cyanide concentration. Panels (a)--(c) use $0$, $10^{-4}$, and $10^{-3}\,\mathrm{M}$, respectively; fixed conditions are nominal metal-ion concentration $10^{-4}\,\mathrm{M}$, $[\ce{NH3}]_{\mathrm{tot}}=0.02\,\mathrm{M}$, and $[\mathrm{Gly}]_{\mathrm{tot}}=0.1\,\mathrm{M}$. The black outline denotes pH 10.5--13.5 and $-2.0$ to $2.3$~V vs RHE (shown on the SHE potential axis).}
    \label{fig:Ti_CN_sensitivity_SI}
\end{figure}

